% Source: https://artofproblemsolving.com/community/c3h1340996_2016_comc_a2
% Stable identity: aops-topic-1340996-q02
From $\pi r^2=16\pi$, the radius is $r=4$ cm. The diagonal of an inscribed square is a diameter of the circle, so it has length $8$.

\begin{asy}
usepackage("amsmath");
size(3.6cm);
defaultpen(fontsize(10pt));
pair A=(4,0), B=(0,4), C=(-4,0), D=(0,-4);
draw(circle((0,0),4),linewidth(0.9)); draw(A--B--C--D--cycle,blue+linewidth(0.9));
draw(C--A,red+dashed); dot((0,0));
label("$8$",(0,-0.25),S,red);
label("\mbox{}",(min(currentpicture,user=true)+max(currentpicture,user=true))/2,black); // keeps the text after the figure black when LaTeX typesets the labels inline
\end{asy}

A square with diagonal $d$ has area $\dfrac{d^2}{2}$, so $S=\dfrac{64}{2}=\boxed{32}$.
